\chapter{Validation expérimentale par scénarios}
\label{chap:validation-scenarios}

\section{Scénario effectivement validé}

La campagne disponible comporte une exécution primaire complète du modèle corrigé. Elle vise à vérifier le comportement combiné d'une rupture initiale de produit fini, d'une absence de GPL et d'un retard déterministe sur la première réception de GPL. L'objectif n'est pas de comparer des politiques, mais de contrôler que la chaîne de reconstitution, les relations de traçabilité et les indicateurs temporels restent cohérents jusqu'au PI.

Le run \Agent{RUN\_1773129600000\_1788264883846} utilise le scénario \texttt{SCENARIO DISTRIBUTION}. Une commande cliente de 20 unités est lancée avec un stock initial de produit fini nul. Le GPL est nul, les bouteilles vides valent 250 unités et les accessoires 300 unités. Le retard fournisseur est actif sur \texttt{GPL\_VRAC}, avec un facteur 2 appliqué à la première réception seulement. Le délai ciblé passe de 6 h nominales à 12 h effectives.

Le Dashboard Global contient aussi 7 h nominales et 9 h effectives comme moyennes des configurations matière. Ces moyennes n'écrasent pas la preuve ciblée du GPL, portée par le manifeste, la trace et les messages AER.

\section{Séquence observée}

\begin{table}[H]
  \centering
  \caption{Chronologie compacte du run de validation VSM.}
  \label{tab:chronologie-run-vsm}
  \begin{tabularx}{\textwidth}{@{}L{0.17\textwidth}L{0.26\textwidth}Y@{}}
    \toprule
    Temps simulé & Objet & Observation vérifiée \\
    \midrule
    T=156 s & \Agent{CMD\_1} & Création et réception de la commande de 20 unités. \\
    T=159 s & Stock et Source & Ouverture de l'attente, création de \Agent{REAPPRO\_1} et lancement de Source. \\
    T=233 s & GPL & Réception de 250 unités, puis disponibilité de la matière. \\
    T=235 s & Make & Début d'exécution de la reconstitution. \\
    T=2 314 s & \Agent{REAPPRO\_1} & Vingtième unité terminée et reconstitution close. \\
    T=2 320 s & Deliver & Début du transport de la commande cliente. \\
    T=2 343 s & \Agent{CMD\_1} & Livraison complète, statut tardif, clôture et exports finaux. \\
    \bottomrule
  \end{tabularx}
\end{table}

La relation directe indique que \Agent{REAPPRO\_1} contribue au service de \Agent{CMD\_1}; la relation inverse indique que la commande dépend de cet ordre interne. La feuille de performance compte séparément une commande cliente et un ordre de stock. À la clôture, les 20 unités sont livrées, le stock fini revient à zéro, le GPL vaut zéro, les bouteilles vides 230 et les accessoires 280.

Cette séquence valide Source avant Make lorsque la matière manque, le crédit du stock fini à la fin de la production, le réveil de la commande et le décrément lors de sa satisfaction. Elle confirme aussi que l'ordre interne reste visible dans la trace tout en étant exclu des dénominateurs clients.

\section{Propagation du retard}

Les six transformations décrites dans la \cref{tab:propagation-retard} sont présentes une fois chacune dans Excel et dans l'ABox. Leur charge utile conserve \texttt{GPL\_VRAC}, \Agent{ACT\_1}, les délais de 6 h et 12 h, le retard de 6 h et l'identifiant \Agent{REAPPRO\_1}.

La propagation prouve qu'un écart local de Source est reformulé jusqu'à Deliver. Elle ne prouve pas que la révision minimise un coût global ni que les six heures supplémentaires expliquent seules le statut tardif. Un run sans retard construit avec le même modèle, les mêmes conditions et une graine contrôlée serait nécessaire pour cette attribution.

\section{Résultats de performance}

Les temps de commande vérifient l'identité additive. Le traitement vaut 18 116,83756052 s, l'attente 46 924,40300189 s et le Lead Time 65 041,24056241 s. Le périmètre ZENER porte 59,38852156 s de traitement, 12,69587399 s d'attente et un PCE estimé de 0,70870898. Les scores RL, RS, AG, CO et AM valent respectivement 7,5, 5,94733671, 4,73462025, 0 et 9,70588235. Le PI vaut 6,11881172 sur 10 avec les profils internes du modèle.

Le run contient aussi la détection d'un goulot sur \Agent{sM1.3.1}, une décision locale \texttt{REBALANCE}, un score AHP de 0,730 et les accusés d'application. Ces faits valident l'exécution de cette boucle locale dans le run. L'AHP ne contribue pas directement au PI et aucun effet de la décision sur le score global ne peut être déduit de cette observation unique.

\section{Ce que le run permet de conclure}

Les artefacts valident le modèle corrigé pour une commande, une reconstitution et une perturbation déterministe. Ils établissent la cohérence croisée du manifeste, des relations ontologiques, des flux, des temps VSM, de Responsiveness et de l'agrégation du PI. La structure ISA-95 exportée contient 190 équipements typés, 71 micro-activités et 71 relations d'affectation, conformément à la garde du modèle.

Le run ne valide pas le fonctionnement nominal, Return, une campagne qualité, une panne, plusieurs commandes partageant un ordre de reconstitution, une politique optimale ou une généralisation aux opérations réelles. Il ne fournit ni comparaison causale avec une référence sans retard, ni réplications permettant d'estimer la variabilité. Ces absences délimitent le résultat expérimental; elles ne sont pas remplacées par les anciens chiffres du document source.
